\documentclass[a4paper, 12pt]{article}
\input{prefix_pension.tex}%
\newcommand{\Set}[2]{\left\{ #1 \;\middle\vert\; #2 \right\}}

\newcommand{\pth}[1]{\mleft(#1\mright)}%

\renewcommand{\P}{\operatorname{P}}
\newcommand{\E}{\operatorname{E}}
\newcommand{\Var}{\operatorname{Var}}
\newcommand{\Cov}{\operatorname{Cov}}
\newcommand{\diam}[1]{\operatorname{diam}\mleft(#1\mright)}
\newcommand{\osc}[2]{\operatorname{osc}\mleft(#1,\, #2\mright)}
\newcommand{\supp}[1]{\operatorname{supp}\mleft(#1\mright)}
\newcommand{\indicator}{\mathbf{1}}
\newcommand{\argmin}{\operatorname{argmin}}
\newcommand{\argmax}{\operatorname{argmax}}
\newcommand{\deq}{\stackrel{d}{=}}

\DeclareMathOperator*{\mtimes}{\text{\raisebox{0.25ex}{\scalebox{0.8}{$\bigtimes$}}}}
\DeclareMathOperator*{\motimes}{\text{\raisebox{0.25ex}{\scalebox{0.8}{$\bigotimes$}}}}

\newcommand{\Ex}[1]{\E\mleft[ #1 \mright]}

\newcommand{\ceil}[1]{\mleft\lceil {#1} \mright\rceil}
\newcommand{\floor}[1]{\mleft\lfloor {#1} \mright\rfloor}
\newcommand{\sgn}{\operatorname{sgn}}
\newcommand{\card}{\operatorname{card}}

\newcommand{\brc}[1]{\left\{ {#1} \right\}}
\newcommand{\sbrc}[1]{\left[ {#1} \right]}
\newcommand{\set}[1]{\brc{#1}}%

\newcommand{\abs}[1]{\left\lvert {#1} \right\rvert}%
\newcommand{\norm}[1]{\left\lVert {#1} \right\rVert}
\newcommand{\spanby}{\operatorname{span}}
\newcommand{\esssup}{\operatorname{ess\,sup}}
\newcommand{\inp}[2]{\left\langle {#1},\, {#2} \right\rangle}

\newcommand{\mto}{\overset{m}{\to}}
\newcommand{\wto}{\overset{w}{\to}}
\newcommand{\wstarto}{\overset{w^*}{\to}}
\newcommand{\asto}{\overset{a.s.}{\to}}
\newcommand{\pto}{\overset{p}{\to}}
\newcommand{\dto}{\overset{d}{\to}}
\newcommand{\dsubset}{\overset{d}{\subset}}


\newcommand{\ds}{\displaystyle}%

\newcommand{\R}{\mathbb{R}}%
\newcommand{\N}{\mathbb{N}}
\newcommand{\Q}{\mathbb{Q}}
\newcommand{\Z}{\mathbb{Z}}
\newcommand{\C}{\mathbb{C}}
\newcommand{\F}{\mathcal{F}}
\newcommand{\A}{\mathcal{A}}
\newcommand{\B}{\mathcal{B}}
\newcommand{\T}{\mathcal{T}}
\newcommand{\M}{\mathcal{M}}
\newcommand{\V}{\mathcal{V}}
\renewcommand{\L}{\mathcal{L}}
\renewcommand{\S}{\mathcal{S}}
\newcommand{\I}{\mathcal{I}}
\renewcommand{\H}{\mathcal{H}}
\newcommand{\D}{\mathcal{D}}
\newcommand{\U}{\mathcal{U}}
\newcommand{\G}{\mathcal{G}}
\newcommand{\J}{\mathcal{J}}

\newcommand{\one}{\mathbf{1}}

\newcommand{\exist}{\exists\:}
\newcommand{\forany}{\forall\:}
\renewcommand{\ae}{\quad\text{a.e.}}
\newcommand{\Res}{\operatorname{Res}}
\newcommand{\cl}{\operatorname{cl}}

\newcommand{\iidsim}{\overset{iid}{\sim}}
\newcommand{\Poisson}{\text{Poisson}}
\newcommand{\Ber}{\text{Ber}}
\newcommand{\Exp}{\text{Exp}}

\newcommand{\restrict}[1]{\left.\hspace{-0.2em}\right|_{#1}}
\newcommand{\eval}[3]{\left.{#1}\right|_{#2}^{#3}}

% derivatives
\makeatletter
\renewcommand\d[1]{\mspace{6mu}\mathrm{d}#1\@ifnextchar\d{\mspace{-3mu}}{}}
\makeatother
\def\at{
  \left.
  \vphantom{\int}
  \right|
}

\newcommand{\od}[2]{\frac{d{#1}}{d{#2}}}
\newcommand{\odd}[2]{\frac{d^2{#1}}{d{#2}^2}}
\newcommand{\pd}[2]{\frac{\partial{#1}}{\partial{#2}}}
\newcommand{\pdd}[2]{\frac{\partial^2{#1}}{\partial{#2}^2}}
\newcommand{\pdpd}[3]{\frac{\partial^2{#1}}{\partial{#2}\partial{#3}}}

\makeatletter
\newcommand*{\rom}[1]{\expandafter\@slowromancap\romannumeral #1@}
\makeatother

\title{Extension to Pension Run}

\author{Kai-Jyun Wang\and Kuan-Ming Chen}

\begin{document} 
\setstretch{1.15}

\maketitle

\section{Introduction}


\section{Institutional Background}

\section{Model}
The time is continuous with an infinite horizon. The 
framework follows \cite{achdou2022}, which is a continuous 
time analogue of the Aiyagari–Bewley–Huggett model 
(\cite{aiyagari1994,huggett1993}). We integrates the model 
with the standard life-cycle model as in \cite{huggett1996}. 
In the model, agents have different believes in the 
possible reforms.  

\subsection{Demographics}
Let $b$ be the fertility rate and $m(s)$ the mortality rate. We assume 
that the individuals cannot live longer than $T$ and $n(t, T) = 0$. Let 
$n(t,s)$ be the unnormalized population density. The population dynamic 
is given by 
\begin{equation}
    \pd{n}{t}(t, s) = -\pd{n}{s}(t, s) - m(s)n(t,s), \quad n(t, 0) = b\int_0^T n(t, s)ds.
\end{equation}
Let $N_t = \int_0^T n(t, s)ds$ be the total population. Detrending 
the population equation by $\hat{n}_t = n_t/N_t$ gives 
\begin{equation}
    \pd{\hat{n}}{t}(t, s) = -\pth{\frac{\dot{N}(t)}{N(t)} + m(s)}\hat{n}(t, s) - \pd{\hat{n}}{s}(t, s), 
    \quad \hat{n}(t, 0) = b. 
\end{equation} 
The population growth rate $\eta_t$ is given by 
\begin{equation}
    \begin{split}
        \frac{\dot{N}_t}{N_t} &= \frac{\int_0^T \pd{n}{t}ds}{\int_0^T n(t, s)ds} 
        = \frac{\int_0^T -\pd{n}{s}(t, s) - m(s)n(t,s)ds}{\int_0^T n(t, s)ds} \\
        &= \frac{-(n(t, T) - n(t, 0)) -\int_0^T m(s)n(t,s)ds}{\int_0^T n(t, s)ds} 
        = b - \int_0^T m(s)\hat{n}(t,s)ds
    \end{split}
\end{equation}
In the steady state, the population grows at a constant rate 
$\eta_t = \frac{\dot{N}(t)}{N(t)}$. The steady state $\hat{n}$ is 
given by 
\begin{equation}
    \bar{n}(s) = b\exp\pth{-\int_0^s \eta + m(x)dx}. 
\end{equation}
The normalization condition $\int_0^T \bar{n}(s)ds = 1$ 
determines $\eta$: 
\begin{equation}
    b\int_0^T \exp\pth{-\int_0^s \eta + m(x)dx}ds = 1.
\end{equation}

\subsection{Household}
\paragraph{Retirees.} 
Given the interest rate $r$, a retiree that dose not 
participate the monthly plan maximizes 
\begin{equation}
    V_t^{\ell} = \max \int_t^T e^{-\int_t^s \rho + m(x)dx}\sbrc{u(c_s) + m(s)b(a_s)}ds
\end{equation}
subject to 
\begin{align}
    \dot{a}_t &= r_ta_t + \tr_t - c_t \\ 
    a_t &\geq 0. 
\end{align}
The utility flow is given by 
\begin{equation}
    u(c) = \frac{c^{1-\gamma}}{1-\gamma}
\end{equation} 
and the bequest motive follows \cite{denardi2004} and 
\cite{french2011}: 
\begin{equation}
    b(a) = \theta_b\frac{(a + \kappa)^{1-\gamma}}{1-\gamma}.
\end{equation}
We assume that the transfer $\tr_t$ is a flow from the bequest, 
which is assumed to be distributed uniformly to all the agents alive. 

For a retiree participating the monthly plan, the consumption 
policy solves 
\begin{equation}
    V_t^{m} = \max \int_t^Te^{-\int_t^s \rho + m(x)dx}\sbrc{u(c_s) + m(s)b(a_s)}ds
\end{equation}
subject to 
\begin{align}
    \dot{a}_t &= r_ta_t + \tr_t + p - c_t \\ 
    a_t &\geq 0. 
\end{align}

\paragraph{Workers.}
After aging $T_w$, the individuals can start to work. 
Given interest rate $r_t$ and wage $w_t$, the households 
consume optimally to maximize
\begin{equation}
    V_t^w = \E_t\sbrc{\int_t^Te^{-\int_t^s \rho + m(x)dx}\sbrc{u(c_s) - \psi(s) + m(s)b(a_s)}ds}
\end{equation}
subject to 
\begin{align}
    \dot{a}_t &= r_ta_t + (y_t-\tau(y_t)) + \tr_t - c_t \\
    a_t &\geq 0 \\ 
    y_t &= w_t\exp(\bar{z}(t) + z_t)
\end{align}
where $\bar{z}(t)$ is the earning profile and $z_t\in\mathcal{Z} 
= \set{z_0,\ldots,z_k}$ follows a continuous time Markov chain with 
intensity matrix $\Lambda$. We take $z_0 = -\infty$ for the unemployment 
state. The initial $z$ at $T_w$ is drawn randomly from the 
invariant distribution associated with $\Lambda$. The 
$\psi$ is the disutility to work. The pension tax $\tau$ 
is of the form 
\begin{equation}
    \tau(y) = \tau_0\min\set{y, \bar{y}}
\end{equation}
which is truncated from above.

We allow the workers to retire at any age. Yet, the pension is 
only available for those retiring within the retirement window 
$[T_r-5, T_r+5]$, where $T_r$ is the standard retiremnt 
age. The benefits $\ell$ and $p$ are determined by 
the retirement age
\begin{equation}
    p = (1+0.04\times (\text{retire age} - T_r))\bar{p}, \quad 
    \ell = (1+0.04\times (\text{retire age} - T_r))\bar{\ell}.
\end{equation}
$\bar{p}$ and $\bar{\ell}$ are the benefits of 
retiring at $T_r$. The retirees cannot return to 
work.

\paragraph{Youngs.}
The youngs are born with zero asset and cannot work before $T_w$. 
They are assumed to be hand to mouth so $c_t = \tr_t$. 

\subsection{Firms}
The production function is given by 
\begin{equation}
    Y_t = AK_t^\alpha L_t^{1-\alpha}.
\end{equation}
Given $r_t, w_t$, the firms solve 
\begin{equation}
    \max_{K_t, L_t}AK_t^\alpha L_t^{1-\alpha} - r_tK_t - w_tL_t
\end{equation}
where $L_t$ is the aggregate effective labor. Taking the FOC and 
detrending by $N_t$ gives 
\begin{equation}
    A\pth{\frac{\hat{K}_t}{\hat{L}_t}}^\alpha = A\pth{\frac{K_t}{L_t}}^\alpha = \frac{w_t}{1-\alpha}, \quad \text{and}\quad
    A\pth{\frac{\hat{K}_t}{\hat{L}_t}}^{\alpha-1} = A\pth{\frac{K_t}{L_t}}^{\alpha-1} = \frac{r_t}{\alpha}.
\end{equation}

\subsection{Pension System}
The government collect pension taxes from the labor income and 
finance the pension payments $\ell$ and $p$. Let pension fund balance 
be $B_t$. The budget constraint is given by 
\begin{equation}
    \begin{split}
        dB_t &= \bigg[ r_t B_t + N_t\int \tau_w(w, z)wzd\mathfrak{M}_t \\ 
        &- N_t\int\ell(s)\mathfrak{M}_t(\text{retiring with lumpsum}) - N_t\int p\mathfrak{M}_t(\text{retired with monthly}) \bigg]dt.
    \end{split}
\end{equation}
Detrending gives
\begin{equation}
    \begin{split}
        d\hat{B}_t &= \bigg[ (r_t-\eta_t)\hat{B}_t + \int \tau_w(w, z)wzd\mathfrak{M}_t \\ 
        &- \ell\mathfrak{M}(\text{retiring with lumpsum}) - p\mathfrak{M}(\text{retired with monthly})\bigg]dt.
    \end{split}
\end{equation}

The reform occurs when the fund balance $\hat{B}_t$ becomes 
zero. There are two possible reforms: one is to let the 
pension system collapse and the other to raise the pension 
tax rate such that the pension system is just solvent in the 
long run. Agents have believes on the two possible reforms. 
Among all the agents, $\pi$ proportion believe that 
the pension system would simply collapse in the future 
and $1-\pi$ believe that the pension tax rate would be 
raised to prevent the system from insolvency. 


\subsection{Equilibrium}
Given the initial distribution $\mathfrak{M}_0(da,ds,dz,dp, d\pi)$,  
pension fund balance $\hat{B}_0$, the recursive equilibrium 
consists of optimal consumption policy path $c_t$, 
the optimal retirement policy, the price paths 
$\set{r_t, w_t}$, the fund balance path $\hat{B}_t$ and the 
distribution $\mathfrak{M}_t$ such that 
\begin{thmenum}
    \item Given the prices $r_t, w_t$ and $\hat{B}_t$, the optimal 
    consumption policies and the retirement policies solves the 
    individuals optimization problem. 
    \item Given the prices $r_t$ and $w_t$, the firms optimize.
    \item The prices are such that the asset market and the 
    labor market clear: 
    \begin{align}
        \hat{K}_t &= \int a d\mathfrak{M}_t + \hat{B}_t \\ 
        \hat{L}_t &= \int \exp(\bar{z}(s) + z)\one\set{\text{working}}d\mathfrak{M}_t. 
    \end{align}
    \item The individuals optimal retirement decision $R_t$ is 
    consistent with the fund balance path $\hat{B}_t$. 
    \item The distribution solves the corresponding 
    Kolmogorov forward equation.
\end{thmenum}

\section{Results}

\section{Conclusion}

\newpage
\section*{\centering Appendix}
\begin{appendices}
\section{Numerical Detail}
\end{appendices}


\bibliographystyle{apacite}
\bibliography{pension}

\end{document}